\scanpage{152}
$n$, trivial pour $n=0$. Cela nous ramène au cas
où $C^1=\ldots=C^{n-1}=0$, donc à prouver le

\uline{Lemme} Il existe des formes linéaires $\alpha^1,\ldots,\alpha^N$
entières en $X^1,\ldots X^r$ telles que
\[
\sigma^1((\alpha^i))=\ldots=\sigma^{n-1}((\unsure{\alpha_i|\alpha^i}))=0,\quad
\sigma^n((\unsure{\alpha_i|\alpha^i}))=C^n\ \text{polynome donné}
\]

\diagram{\begin{tikzcd}[ampersand replacement=\&,row sep=1.6em,column sep=large,execute at end picture={\draw (current bounding box.north west) -- (current bounding box.north east) -- (current bounding box.south east);}]
\mathrm Z[\unsure{\xi_1,\ldots,\xi_n|\zeta_1,\ldots,\zeta_n|\xi_1,\ldots}] \& {}\\
E_T\arrow[d] \& {}\\
B_T \& {}\\
\mathbb Z[\widehat T]\arrow[r,"\unsure{\sim}"] \& H^*(B_T,\mathbb Z)\\
{} \& A(T)\arrow[ul,"\text{alg}"]\arrow[u,"\text{top via Chern}"']\\
\mathbb Z[\widehat T]\arrow[r,"{\text{\unsure{hom. car.}}}"{draw,circle,inner sep=1pt,above}] \& H^*(X,\mathbb Z)\\
{} \& A(T)\arrow[ul]\arrow[u,"\text{top. via Chern}"']
\end{tikzcd}}

Il faudrait prouver que
$\mathbb Z[\widehat T]$ est l'anneau de
Chern attaché à $A(T)$
$\simeq\mathbb Z(\widehat T)$
% archive sheet number: 151
% printed form (vertical paper watermark)
\tikz[baseline=-.4ex,scale=.15]{\foreach \a in {0,45,...,315}{\draw (\a:.52) circle (.22);} \draw (0,0) circle (.21);} ARGENTINE \tikz[baseline=-.2ex,scale=.17]{\draw (0,0)--(1.6,0)--(1.8,.33)--(.5,.33)--(.32,.13)--(.92,.13);} TYPEWRITER \tikz[baseline=-.4ex,scale=.15]{\foreach \a in {0,45,...,315}{\draw (\a:.52) circle (.22);} \draw (0,0) circle (.21);} \unsure{AR}
